\documentclass[9pt,aspectratio=169,xcolor=table]{beamer}

%~~~~~~~~~~~~~~~~~~~~~~~~~~~~~~~~~~~~~~~~~~~~~~~~~~~~~~~~~~~~~~~~~~~~~~~~~~~~~~
\usepackage[sfdefault]{roboto}
\usepackage[utf8]{inputenc}
\usepackage[T1]{fontenc}
%~~~~~~~~~~~~~~~~~~~~~~~~~~~~~~~~~~~~~~~~~~~~~~~~~~~~~~~~~~~~~~~~~~~~~~~~~~~~~~

\usepackage{styles/fluxmacros}
\usefolder{styles}
\usetheme[style=red]{flux}

%~~~~~~~~~~~~~~~~~~~~~~~~~~~~~~~~~~~~~~~~~~~~~~~~~~~~~~~~~~~~~~~~~~~~~~~~~~~~~~
% SKEE1033 maroon palette override (see ecp01.tex for rationale).
\definecolor{primaryLight}{HTML}{9B2242}
\definecolor{primary}{HTML}{7B1E3A}
\definecolor{secondary}{HTML}{3A5A6B}
\definecolor{tertiary}{HTML}{8C6A3F}
%~~~~~~~~~~~~~~~~~~~~~~~~~~~~~~~~~~~~~~~~~~~~~~~~~~~~~~~~~~~~~~~~~~~~~~~~~~~~~~

\usepackage{booktabs}
\usepackage{colortbl}
\usepackage{ragged2e}
\usepackage{amssymb}
\usepackage{multirow}
\usepackage{tikz}
\usetikzlibrary{arrows.meta, positioning, fit, backgrounds, calc, shapes.geometric}
\setbeamertemplate{itemize items}[square]
\usepackage[scaled=.92]{beramono}
\usepackage{listings}
\usepackage{array}
\definecolor{stripe}{gray}{0.94}
\newcommand{\striped}{\rowcolors{1}{white}{stripe}}

%~~~~~~~~~~~~~~~~~~~~~~~~~~~~~~~~~~~~~~~~~~~~~~~~~~~~~~~~~~~~~~~~~~~~~~~~~~~~~~
% Code listing styles. Python is the lstlisting default via \lstset below.
\definecolor{codebg}{HTML}{FBF2F4}
\definecolor{codegreen}{rgb}{0,0.45,0}
\definecolor{codestring}{rgb}{0.58,0.1,0.2}
\lstdefinestyle{skpython}{
    basicstyle=\ttfamily\footnotesize,
    keywordstyle=\color{primary}\bfseries,
    commentstyle=\itshape\color{codegreen},
    stringstyle=\ttfamily\color{codestring},
    backgroundcolor=\color{codebg},
    breaklines=true,
    keepspaces=true,
    showspaces=false,
    showstringspaces=false,
    frame=single,
    rulecolor=\color{primary!40},
    numbers=none,
}
\lstdefinestyle{skoutput}{
    basicstyle=\ttfamily\footnotesize,
    backgroundcolor=\color{secondary!8},
    breaklines=true,
    keepspaces=true,
    showspaces=false,
    showstringspaces=false,
    frame=single,
    rulecolor=\color{secondary!50},
    numbers=none,
}
\lstset{language=Python, style=skpython}
%~~~~~~~~~~~~~~~~~~~~~~~~~~~~~~~~~~~~~~~~~~~~~~~~~~~~~~~~~~~~~~~~~~~~~~~~~~~~~~

\definecolor{cycfill}{HTML}{F3DCE4}
\definecolor{cycline}{HTML}{7B1E3A}
\definecolor{cycdofill}{HTML}{E8EEF0}
\definecolor{cycdoline}{HTML}{3A5A6B}
\tikzset{
  cycstep/.style = {font=\sffamily\small, align=center, rounded corners=2pt,
                     draw=cycline, fill=cycfill, line width=0.7pt,
                     minimum height=7mm, minimum width=20mm, inner sep=3pt},
  cycdo/.style   = {cycstep, fill=cycdofill, draw=cycdoline},
  flow/.style    = {-{Stealth[length=2mm]}, thick, draw=cycline},
}

%%%%%%%%%%%%%%%%%%%%%%%%%%%%%%%%%%%
%% SECTION SEPARATOR (ToC recap) %%
%%%%%%%%%%%%%%%%%%%%%%%%%%%%%%%%%%%
\AtBeginSection[]
{
{
    \setbeamercolor{background canvas}{bg=primaryLight!6}
    \begin{frame}[plain]
        \frametitle{Table of Contents}
        \tableofcontents[currentsection]
    \end{frame}
}
}
%%%%%%%%%%%%%%%%%%%%%%%%%%%%%%%%%%%

\title{Numbers, Variables and Units}
\subtitle{\emph{Engineering Computation with Python} --- Chapter 3}
\author{Mun\'{i}m Zabidi}
\institute{\color{primary}Faculty of Electrical Engineering, Universiti Teknologi Malaysia}
\date{\today}
\titlegraphic{assets/logoutmwhite.png}

\begin{document}

\titlepage

\begin{frame}{Table of Contents}
    \tableofcontents
\end{frame}

%===============================================================================
\section{Beyond Simple Numbers}
%===============================================================================

\begin{frame}{Engineering Data Is Often More Complex Than This}
    \leavevmode
    Chapter 2 covered values that hold a single number or word.
    Engineering data is often more complex:

    \vspace{2mm}
    \begin{itemize}
        \item An AC voltage needs \alert{two components} at once.
        \item A measurement run produces a \alert{whole series} of
              readings.
        \item Every quantity carries a \alert{unit} that decides whether
              the answer is right.
    \end{itemize}

    \vspace{3mm}
    \begin{block}{This chapter}
        Richer types for all three --- complex numbers, strings, and
        four collection types --- then how to represent quantities like
        $4.7\,\text{k}\Omega$ and $10\,\mu\text{F}$ correctly.
    \end{block}
\end{frame}

\begin{frame}[fragile]{Complex Numbers in Python}
    \leavevmode
    Python natively supports complex numbers: $a + bj$, where $j$ is the
    imaginary unit ($j^2=-1$).

\begin{lstlisting}
c1 = 3 + 4j   # real part 3, imaginary part 4
c2 = -2 - 5j
print(c1)          # (3+4j)
print(type(c1))    # <class 'complex'>

c = complex(3, 4)  # equivalent to 3 + 4j
\end{lstlisting}

    \vspace{1mm}
    {\small Python uses \texttt{j}, not the mathematician's \texttt{i}
    --- matching electrical-engineering convention, where $i$ already
    means current.}
\end{frame}

\begin{frame}[fragile]{Working with Complex Numbers}
    \leavevmode
    \begin{columns}[t]
        \begin{column}{.5\textwidth}
\begin{lstlisting}
c = 3 + 4j
print(c.real)  # 3.0
print(c.imag)  # 4.0

c1 = 3 + 4j
c2 = 1 - 2j
print(c1 + c2)  # (4+2j)
print(c1 * c2)  # (11-2j)
\end{lstlisting}
        \end{column}
        \begin{column}{.5\textwidth}
\begin{lstlisting}
c = 3 + 4j
print(c.conjugate())
# (3-4j)

print(abs(c))
# 5.0  (magnitude)
\end{lstlisting}
            \vspace{2mm}
            {\small The four arithmetic operators work exactly as on
            real numbers --- Python handles the real/imaginary parts for
            you.}
        \end{column}
    \end{columns}
\end{frame}

\begin{frame}{Why Complex Numbers Matter in EE}
    \leavevmode
    \begin{block}{Signal processing and electrical engineering}
        \small Commonly used for oscillations or waves --- AC circuit
        analysis represents voltage and current as complex phasors,
        making impedance calculations straightforward.
    \end{block}

    \vspace{3mm}
    \begin{block}{Mathematics}
        \small Useful for solving equations involving square roots of
        negative numbers.
    \end{block}

    \vspace{3mm}
    {\small Complex numbers for AC circuits return in full in Chapter 8.}
\end{frame}

\begin{frame}[fragile]{Strings: Text as a Sequence}
    \leavevmode
    A \alert{string} is a sequence of characters enclosed in quotes.
    Strings are \alert{immutable} --- they cannot be changed after
    creation.

\begin{lstlisting}
string1 = 'Hello'
string2 = "World"
string3 = '''A multi-line
string.'''
\end{lstlisting}

    \vspace{1mm}
    {\small Single and double quotes are interchangeable; triple quotes
    span several lines.}
\end{frame}

\begin{frame}[fragile]{Indexing and Slicing Strings}
    \leavevmode
\begin{lstlisting}
text = "Python"
first_char = text[0]   # 'P'
last_char = text[-1]   # 'n'

text = "Python Programming"
slice1 = text[0:6]  # 'Python'
slice2 = text[7:]   # 'Programming'
\end{lstlisting}

    \vspace{2mm}
    \begin{block}{Remember}
        Counting starts from zero. A negative index counts back from
        the end. \texttt{stop} in a slice is \alert{excluded}.
    \end{block}
\end{frame}

\begin{frame}[fragile]{String Methods}
    \leavevmode
    Because strings are immutable, every method below returns a
    \alert{new} string rather than altering the original.

\begin{lstlisting}
text = " Python Programming "

text.upper()     # ' PYTHON PROGRAMMING '
text.lower()     # ' python programming '
text.strip()     # 'Python Programming'
text.replace("Python", "Java")
# ' Java Programming '
text.split()     # ['Python', 'Programming']
\end{lstlisting}
\end{frame}

%===============================================================================
\section{Collections}
%===============================================================================

\begin{frame}{Four Ways to Hold Several Values}
    \leavevmode
    So far every variable has held exactly one value. Engineering work
    rarely stays that simple: a measurement run produces a series of
    readings, a circuit has a set of component values.

    \vspace{3mm}
    \begin{columns}[t]
        \begin{column}{.25\textwidth}
            \begin{block}{\small List}
                \scriptsize Ordered, changeable
            \end{block}
        \end{column}
        \begin{column}{.25\textwidth}
            \begin{block}{\small Tuple}
                \scriptsize Ordered, fixed
            \end{block}
        \end{column}
        \begin{column}{.25\textwidth}
            \begin{block}{\small Set}
                \scriptsize Unordered, unique
            \end{block}
        \end{column}
        \begin{column}{.25\textwidth}
            \begin{block}{\small Dictionary}
                \scriptsize Looked up by name
            \end{block}
        \end{column}
    \end{columns}

    \vspace{3mm}
    {\small Choosing the wrong one is a common source of awkward code.}
\end{frame}

\begin{frame}[fragile]{Lists: The Everyday Workhorse}
    \leavevmode
\begin{lstlisting}
readings = [26.5, 26.6, 26.5, 26.7, 26.6]
print(readings[0])     # 26.5
print(readings[-1])    # 26.6
print(len(readings))   # 5

readings.append(27.0)  # add to the end
readings[1] = 99.9     # lists are mutable
\end{lstlisting}

    \vspace{2mm}
    \begin{block}{Ordered and changeable}
        A list keeps items in the order you put them, and you can add,
        remove, or reassign elements after creation.
    \end{block}
\end{frame}

\begin{frame}[fragile]{Slicing a List}
    \leavevmode
\begin{lstlisting}
data = [10, 20, 30, 40, 50, 60]

print(data[1:4])    # [20, 30, 40]
print(data[:3])     # [10, 20, 30]
print(data[3:])     # [40, 50, 60]
print(data[::2])    # [10, 30, 50]  -- every 2nd
print(data[::-1])   # reversed
\end{lstlisting}

    \vspace{2mm}
    {\small \texttt{[start:stop:step]} --- extract a window, a tail, or
    every $n$th element. \texttt{stop} is always excluded.}
\end{frame}

\begin{frame}[fragile]{Tuples: Values That Belong Together}
    \leavevmode
\begin{lstlisting}
point = (3, 4)
rgb_color = (255, 128, 0)

red, green, blue = (255, 128, 0)  # unpacking
print(red, green, blue)  # 255 128 0

# point[0] = 5  # TypeError -- tuples are immutable
\end{lstlisting}

    \vspace{2mm}
    \begin{alertblock}{Immutable by design}
        The right choice when values belong together permanently and
        should not be edited by accident: an $(x,y)$ coordinate, an RGB
        colour, a fixed pair of circuit terminals.
    \end{alertblock}
\end{frame}

\begin{frame}{Tuples Reappear: NumPy Shapes}
    \leavevmode
    You will meet tuples again sooner than you might expect.

    \vspace{3mm}
    \begin{block}{The shape of a NumPy array is a tuple}
        \texttt{np.zeros((3, 3))} --- the dimensions are passed inside
        an \alert{extra pair of brackets}, because the shape itself is
        a single tuple argument.
    \end{block}

    \vspace{3mm}
    {\small This becomes essential terminology from Chapter 6 onward.}
\end{frame}

\begin{frame}[fragile]{Sets: Only Distinct Values Matter}
    \leavevmode
    A \alert{set} is an unordered collection that automatically discards
    duplicates.

\begin{lstlisting}
readings = [26.5, 26.6, 26.5, 26.7, 26.6]
unique = set(readings)
print(sorted(unique))   # [26.5, 26.6, 26.7]
print(len(unique))      # 3 distinct values

a = {1, 2, 3}
b = {3, 4}
print(a | b)   # union: {1, 2, 3, 4}
print(a & b)   # intersection: {3}
\end{lstlisting}

    \vspace{1mm}
    {\small No order means no indexing --- \texttt{unique[0]} is an
    error.}
\end{frame}

\begin{frame}[fragile]{Dictionaries: Looked Up by Name}
    \leavevmode
    A \alert{dictionary} stores values looked up by a \alert{key} rather
    than a numeric position --- \texttt{key: value} pairs.

\begin{lstlisting}
components = {"R1": 4700, "R2": 100, "C1": 100e-6}

print(components["R1"])     # 4700
components["R3"] = 220      # add a new entry

for name, value in components.items():
    print(f"{name}: {value}")
\end{lstlisting}

    \vspace{1mm}
    {\small A list would force you to remember that position 0 means
    R1 --- a fact nothing in the code records.}
\end{frame}

\begin{frame}{Choosing a Collection}
    \leavevmode
    \small
    \striped
    \begin{center}
    \begin{tabular}{p{1.4cm}p{2.5cm}p{1.3cm}p{1.3cm}p{4.3cm}}
        \toprule
        \textbf{Type} & \textbf{Written} & \textbf{Ordered?} & \textbf{Mutable?} & \textbf{Use it when} \\
        \midrule
        list  & \texttt{[1, 2, 3]}   & yes & yes & values form a changing sequence \\
        tuple & \texttt{(1, 2, 3)}   & yes & no  & values belong together, fixed \\
        set   & \texttt{\{1, 2, 3\}} & no  & yes & only distinct values matter \\
        dict  & \texttt{\{"a": 1\}}  & yes & yes & values are looked up by name \\
        \bottomrule
    \end{tabular}
    \end{center}

    \vspace{2mm}
    \begin{block}{For numerical work}
        The \alert{list} is by far the most important --- and it is
        itself soon superseded by the NumPy \alert{array} in Chapter 6.
    \end{block}
\end{frame}

%===============================================================================
\section{Mutability}
%===============================================================================

\begin{frame}[fragile]{Mutable vs. Immutable}
    \leavevmode
    \begin{columns}[t]
        \begin{column}{.5\textwidth}
            \begin{block}{Mutable: can change in place}
                \scriptsize list, set, dictionary
            \end{block}
\begin{lstlisting}
a = [1, 2, 3]
b = a          # b refers to
               # the SAME list
b.append(4)
print(a)
# [1, 2, 3, 4]  -- a changed too!
\end{lstlisting}
        \end{column}
        \begin{column}{.5\textwidth}
            \begin{block}{Immutable: cannot change}
                \scriptsize int, float, str, tuple, complex, bool
            \end{block}
\begin{lstlisting}
x = 5
y = x
y = y + 1
print(x)
# 5  -- x is unaffected
\end{lstlisting}
        \end{column}
    \end{columns}
\end{frame}

\begin{frame}{Why Mutability Matters}
    \leavevmode
    \begin{alertblock}{The hazard}
        Assigning a mutable object to a new name does \alert{not} copy
        it --- both names point to the same underlying data. A change
        through one name is visible through the other.
    \end{alertblock}

    \vspace{3mm}
    \begin{block}{Practical consequence}
        \small Passing a list into a function and modifying it inside
        the function changes the \emph{caller's} list too, unless you
        deliberately make a copy first. This surprises nearly every
        beginner at least once.
    \end{block}
\end{frame}

%===============================================================================
\section{Engineering Units}
%===============================================================================

\begin{frame}{Every Number Carries a Unit}
    \leavevmode
    A resistance of $470$ is meaningless; a resistance of $470\,\Omega$
    is a component you can hold.

    \vspace{3mm}
    \begin{alertblock}{Python does not track units}
        It will happily add a voltage to a resistance and give you a
        number. Keeping track of units is \alert{your} responsibility
        --- one of the most common sources of error for beginners.
    \end{alertblock}
\end{frame}

\begin{frame}{Quantities You Will Meet Most}
    \leavevmode
    \small
    \striped
    \begin{center}
    \begin{tabular}{p{2.6cm}p{1.1cm}p{1.8cm}p{4.0cm}}
        \toprule
        \textbf{Quantity} & \textbf{Sym.} & \textbf{Unit} & \textbf{Relationship} \\
        \midrule
        Voltage      & $V$ & volt (V)    & $V = IR$ \\
        Current      & $I$ & ampere (A)  & $I = V/R$ \\
        Resistance   & $R$ & ohm ($\Omega$) & $R = V/I$ \\
        Power        & $P$ & watt (W)    & $P = VI$ \\
        Capacitance  & $C$ & farad (F)   & $i = C\,dv/dt$ \\
        Frequency    & $f$ & hertz (Hz)  & $f = 1/T$ \\
        \bottomrule
    \end{tabular}
    \end{center}
\end{frame}

\begin{frame}{SI Prefixes You Will Use Constantly}
    \leavevmode
    \small
    A capacitor might be $10$ microfarads while a resistor is
    $4.7$ kilohms --- writing either in full decimal form is error-prone.

    \vspace{2mm}
    \striped
    \begin{center}
    \begin{tabular}{llll}
        \toprule
        \textbf{Prefix} & \textbf{Sym.} & \textbf{Mult.} & \textbf{Example} \\
        \midrule
        kilo  & k      & $10^{3}$   & $4.7\,\text{k}\Omega$ \\
        milli & m      & $10^{-3}$  & $20\,\text{mA}$ \\
        micro & $\mu$  & $10^{-6}$  & $10\,\mu\text{F}$ \\
        nano  & n      & $10^{-9}$  & $5\,\text{ns}$ \\
        pico  & p      & $10^{-12}$ & $22\,\text{pF}$ \\
        \bottomrule
    \end{tabular}
    \end{center}
\end{frame}

\begin{frame}[fragile]{Scientific Notation in Python}
    \leavevmode
    Python has no notion of ``kilo'' or ``micro'' --- prefixed values
    must be entered as plain numbers, using \alert{scientific notation}.

\begin{lstlisting}
R = 4.7e3      # 4.7 kilohms = 4700 ohms
C = 100e-6     # 100 microfarads = 0.0001 farads
I = 20e-3      # 20 milliamps = 0.02 amps
\end{lstlisting}

    \vspace{2mm}
    {\small The letter \texttt{e} means ``times ten to the power of.''
    This is the clearest way to write a prefixed engineering value.}
\end{frame}

\begin{frame}{Working Consistently in Base Units}
    \leavevmode
    \begin{block}{The rule}
        Convert every quantity to its \alert{base SI unit} (ohms, farads,
        seconds, amps) before calculating, then convert the \emph{result}
        back to a convenient prefix for display.
    \end{block}

    \vspace{3mm}
    \begin{alertblock}{Worked check: RC time constant}
        $R = 4.7\,\text{k}\Omega = 4700\,\Omega$, $C = 100\,\mu\text{F} =
        0.0001\,\text{F}$:
        \[
        \tau = RC = 4700 \times 0.0001 = 0.47\,\text{s}.
        \]
        Units confirm the result: ohms $\times$ farads $=$ seconds.
    \end{alertblock}
\end{frame}

%===============================================================================
\section{Conventions and Constants}
%===============================================================================

\begin{frame}{Naming Conventions: PEP 8}
    \leavevmode
    Python does not care what you call things --- \texttt{supply\_voltage},
    \texttt{SupplyVoltage}, and \texttt{sv} are all legal. Other engineers
    \alert{do} care.

    \vspace{3mm}
    \begin{columns}[t]
        \begin{column}{.33\textwidth}
            \begin{block}{\scriptsize snake\_case}
                \scriptsize variables, functions \\
                \texttt{\scriptsize supply\_voltage}
            \end{block}
        \end{column}
        \begin{column}{.33\textwidth}
            \begin{block}{\scriptsize SCREAMING\_SNAKE}
                \scriptsize constants \\
                \texttt{\scriptsize MAX\_CURRENT}
            \end{block}
        \end{column}
        \begin{column}{.33\textwidth}
            \begin{block}{\scriptsize CapWords}
                \scriptsize classes \\
                \texttt{\scriptsize Resistor}
            \end{block}
        \end{column}
    \end{columns}

    \vspace{3mm}
    {\small The style you will use constantly is \alert{snake\_case}.}
\end{frame}

\begin{frame}[fragile]{Naming Well: A Practical Example}
    \leavevmode
\begin{lstlisting}
# Conventional -- every name announces what it is
MAX_CURRENT = 0.02              # a constant
supply_voltage = 5.0            # a variable

def series_resistor(v_supply, v_led, current):
    return (v_supply - v_led) / current
\end{lstlisting}

    \vspace{2mm}
    \begin{block}{Two further habits}
        Name the \emph{quantity}, not its type (\texttt{voltage} beats
        \texttt{voltage\_float}). Reserve single letters for symbols
        that match the mathematics (\texttt{R}, \texttt{C}, \texttt{f},
        \texttt{t}).
    \end{block}
\end{frame}

\begin{frame}[fragile]{Constants}
    \leavevmode
    Some quantities never change: the speed of light, a supply rail
    fixed by hardware. Python has no \texttt{const} keyword --- it uses
    the \texttt{SCREAMING\_SNAKE\_CASE} convention instead.

\begin{lstlisting}
PI = 3.14159
GRAVITY = 9.81
SPEED_OF_LIGHT = 299792458
\end{lstlisting}

    \vspace{2mm}
    \begin{alertblock}{The capitals carry no authority}
        Nothing stops a later line from reassigning \texttt{GRAVITY} ---
        a constant in Python is a \alert{promise between programmers},
        not a guarantee from the language.
    \end{alertblock}
\end{frame}

%===============================================================================
\section{Summary}
%===============================================================================

\begin{frame}{What You Should Take Away}
    \leavevmode
    \small
    \begin{itemize}
        \item \textbf{Complex numbers} (\texttt{3+4j}) and
              \textbf{strings} are Python's richer scalar types, both
              with their own arithmetic and methods.
        \item Four \textbf{collections}: list (ordered, mutable), tuple
              (ordered, fixed), set (unordered, unique), dictionary
              (keyed lookup).
        \item \textbf{Mutability} matters: assigning a mutable object to
              a new name shares the same underlying data.
        \item Every number carries a \textbf{unit}. Work in \textbf{base
              SI units}, using scientific notation (\texttt{4.7e3}) for
              prefixed values.
    \end{itemize}
\end{frame}

\begin{frame}{What You Should Take Away (continued)}
    \leavevmode
    \small
    \begin{itemize}
        \item \textbf{PEP 8} naming: \texttt{snake\_case} for variables
              and functions (your everyday style), \texttt{SCREAMING\_}
              \texttt{SNAKE\_CASE} for constants.
        \item A Python \textbf{constant} is a naming convention, not an
              enforced guarantee --- honouring it is what makes it
              useful.
        \item The \textbf{list} is the most important collection for
              numerical work --- and is superseded by the NumPy
              \textbf{array} from Chapter 6 onward.
    \end{itemize}

    \vspace{4mm}
    \begin{center}
        \large\color{primary}\textbf{Next: Decisions and Repetition}
    \end{center}
\end{frame}

\end{document}
